\subsection{Sensitivity to the accepted input information}\label{sec:sourcepolicy}
The conservative choice depends on which submitted observations are accepted.
Of the 112 retained point entries, 64 are labelled 2026, two 2025 and 46 earlier;
28 carry non-\texttt{ok} source flags: 16 \texttt{MODELLED}, eight \texttt{STALE},
three \texttt{SUSPECT} and one \texttt{NATIONAL}.
For a controlled test, replace selected point utilities by $[0,1]$, keeping
weights, all supplements and the 48 original non-point bounds unchanged.
The groups are points dated before 2020, after 2025, non-\texttt{ok} points and
their union. The union counts overlapping entries once.

\begin{table}[!ht]\centering
\caption{Recommendation after withdrawing selected submitted point information. Width is the mean baseline-weight score range on the 0--100 scale.}\label{tab:vintage}
\begin{tabular}{l r l l r}\toprule
Withdrawn group & Cells & Baseline & Equal & Width\\\midrule
Retained input & 0 & Zurich & Zurich & 16.06 \\
Before 2020 & 13 & Singapore & Zurich & 21.14 \\
After 2025 & 64 & Zurich & Zurich & 48.42 \\
Non-ok flags & 28 & Berlin & Paris & 29.11 \\
Union & 95 & Singapore & Singapore & 68.79 \\
\bottomrule\end{tabular}\end{table}


Withdrawing the 28 flagged points changes the recommendation to Berlin under
baseline weights and Paris under equal weights. Removing the 13 pre-2020 points
selects Singapore under baseline weights; equal weights retain Zurich.
Withdrawing the 64 later-dated points retains Zurich under both references,
while widening the baseline mean range from 16.06 to 48.42. The 95-entry union
selects Singapore under both. All sixteen cities remain possible champions
in every case. These withdrawals test dependence on accepted information:
the labels and date cutoffs are not calibrated error probabilities or a
judgment that every flagged entry is wrong. The experiment preserves the
primary data and does not reconstruct a same-year-only matrix.

Withholding Zurich's homicide-rate point alone selects Singapore under baseline
weights and Oslo under equal weights. Withholding Amsterdam's alone retains
Zurich under both. Each single-point test keeps the other 159 entries fixed;
no official count is inserted as an assumed replacement rate.

Holding the other 159 entries fixed while varying this point rate $x_Z$, Zurich ties Singapore in maximum regret at about $x_Z=0.23002$ per 100,000 under baseline weights, and Oslo at about $x_Z=0.90473$ under equal weights. Below each respective cutoff Zurich uniquely minimizes regret; above it Zurich does not. These are hypothetical point-rate cutoffs, not observed rates or confidence limits.
